\documentclass[a4paper,12pt]{article}
\usepackage{amsmath, amssymb}
\usepackage{xcolor}
\usepackage[most]{tcolorbox}
\usepackage{multicol}
\usepackage{geometry}
\usepackage{graphicx}
\usepackage{hyperref}

\geometry{left=1.5cm, right=1.5cm, top=2cm, bottom=2cm}

\author{}
\date{}

% Style de boîte pour les exercices
\tcbset{
    colframe=purple!70,
    colback=white,
    coltitle=black,
    fonttitle=\bfseries,
    sharp corners,
    boxrule=0.5mm,
    colbacktitle=purple!20!white,
    center title,
    boxsep=4pt,
}

% Style de boîte pour les solutions
\newtcolorbox{solutionbox}[1][]{
    colframe=green!70,
    colback=white,
    coltitle=black,
    fonttitle=\bfseries,
    sharp corners,
    boxrule=0.5mm,
    colbacktitle=green!20!white,
    title=Solution,
    boxsep=4pt,
    breakable,
    #1
}

\begin{document}

% Exercice 10
\begin{tcolorbox}[title=Exercice 10 - {[}Calculer{]}]
Soit \( n \geq 2 \). Résoudre l’équation suivante dans \( \mathbb{C} \) :
\[
z^n = \bar{z}.
\]
\end{tcolorbox}

\begin{solutionbox}
Soit \( z \in \mathbb{C} \) tel que \( z^n = \bar{z} \).

Écrivons \( z \) sous forme exponentielle :
\[
z = r e^{i\theta}, \quad r \geq 0, \theta \in \mathbb{R}.
\]
Alors, son conjugué est :
\[
\bar{z} = r e^{-i\theta}.
\]

L'équation devient :
\[
(z)^n = \bar{z} \quad \Rightarrow \quad (r e^{i\theta})^n = r e^{-i\theta}.
\]

Simplifions :
\[
r^n e^{in\theta} = r e^{-i\theta}.
\]

Divisons les deux membres par \( r \) (supposé non nul) :
\[
r^{n-1} e^{in\theta} = e^{-i\theta}.
\]

Égalité des modules :
\[
r^{n-1} = 1 \quad \Rightarrow \quad r = 1.
\]

Égalité des arguments :
\[
n\theta = -\theta + 2k\pi, \quad k \in \mathbb{Z}.
\]

Cela donne :
\[
n\theta + \theta = 2k\pi \quad \Rightarrow \quad (n + 1)\theta = 2k\pi.
\]

Donc :
\[
\theta = \dfrac{2k\pi}{n+1}, \quad k \in \mathbb{Z}.
\]

Comme \( \theta \) est défini modulo \( 2\pi \), les valeurs distinctes de \( \theta \) sont :
\[
\theta_k = \dfrac{2k\pi}{n+1}, \quad k = 0, 1, 2, \dots, n.
\]

Ainsi, les solutions sont :
\[
z_k = e^{i\theta_k} = e^{i\frac{2k\pi}{n+1}}, \quad k = 0, 1, 2, \dots, n.
\]

\textbf{Conclusion} : Les solutions de l'équation sont les racines \((n+1)\)-ièmes de l'unité.
\end{solutionbox}

% Exercice 11
\begin{tcolorbox}[title=Exercice 11 - {[}Calculer{]}]
Soit \( n \geq 2 \). Calculer le produit des racines \( n^{\text{e}} \) de l’unité.
\end{tcolorbox}

\begin{solutionbox}
Les racines \( n \)-ièmes de l'unité sont les nombres complexes de la forme :
\[
z_k = e^{i\frac{2k\pi}{n}}, \quad k = 0, 1, 2, \dots, n-1.
\]

Le produit de ces racines est :
\[
P = \prod_{k=0}^{n-1} z_k = \prod_{k=0}^{n-1} e^{i\frac{2k\pi}{n}} = e^{i\frac{2\pi}{n} \sum_{k=0}^{n-1} k}.
\]

Calculons la somme :
\[
\sum_{k=0}^{n-1} k = \dfrac{(n-1)n}{2}.
\]

Donc :
\[
P = e^{i\frac{2\pi}{n} \cdot \dfrac{(n-1)n}{2}} = e^{i\pi(n-1)}.
\]

Comme \( e^{i\pi(n-1)} = (-1)^{n-1} \), on obtient :
\[
P = (-1)^{n-1}.
\]

\textbf{Conclusion} : Le produit des racines \( n \)-ièmes de l'unité est \( (-1)^{n-1} \).
\end{solutionbox}

% Exercice 12
\begin{tcolorbox}[title={Exercice 12 - {[}Calculer, Chercher{]}}]
Déterminer tous les nombres complexes \( z \) tels que les points d’affixes \( z \), \( z^2 \) et \( z^4 \) soient alignés.
\end{tcolorbox}

\begin{solutionbox}
Pour que les points d'affixes \( z \), \( z^2 \) et \( z^4 \) soient alignés, il faut que les vecteurs \( \overrightarrow{z z^2} \) et \( \overrightarrow{z z^4} \) soient colinéaires.

Les vecteurs \( \overrightarrow{z z^2} \) et \( \overrightarrow{z z^4} \) sont colinéaires si, et seulement si, les nombres complexes \( z^2 - z \) et \( z^4 - z \) sont colinéaires, c'est-à-dire s'il existe un réel \( \mu \) tel que :
\[
\dfrac{z^2 - z}{z^4 - z} \in \mathbb{R}
\]
Calculons ce quotient :
\[
\dfrac{z^2 - z}{z^4 - z} = \dfrac{z(z - 1)}{z(z^3 - 1)} = \dfrac{z - 1}{z^3 - 1}
\]
Donc, la condition est que le nombre complexe :
\[
Q = \dfrac{z - 1}{z^3 - 1}
\]
soit réel.

Notons que :
\[
z^3 - 1 = (z - 1)(z^2 + z + 1)
\]
Ainsi :
\[
Q = \dfrac{z - 1}{(z - 1)(z^2 + z + 1)} = \dfrac{1}{z^2 + z + 1}
\]
Donc :
\[
Q = \dfrac{1}{z^2 + z + 1}
\]
Ainsi, \( Q \) est réel si, et seulement si, \( z^2 + z + 1 \) est réel et \( Q \in \mathbb{R} \).

Donc, il faut que \( z^2 + z + 1 \) soit réel.

\textbf{Soit \( z = x + i y \)} avec \( x, y \in \mathbb{R} \).

Calculons \( z^2 + z + 1 \) :
\[
z^2 + z + 1 = (x + i y)^2 + (x + i y) + 1
\]
Développons :
\[
(x^2 - y^2 + 2i x y) + (x + i y) + 1 = (x^2 - y^2 + x + 1) + i (2 x y + y)
\]
Pour que \( z^2 + z + 1 \) soit réel, il faut que la partie imaginaire soit nulle :
\[
2 x y + y = y (2 x + 1) = 0
\]
Deux cas possibles :

\textbf{Cas 1 :} \( y = 0 \)

Donc \( z \in \mathbb{R} \)

\textbf{Cas 2 :} \( 2 x + 1 = 0 \Rightarrow x = -\dfrac{1}{2} \)

Donc \( z = -\dfrac{1}{2} + i y \)

Alors, \( z^2 + z + 1 \) est réel.

Vérifions que \( Q \) est réel dans ce cas.

Calculons \( z^2 + z + 1 \) :

Pour \( x = -\dfrac{1}{2} \), on a :
\[
x^2 - y^2 + x + 1 = \left( \dfrac{1}{4} - y^2 \right) - \dfrac{1}{2} + 1 = \dfrac{3}{4} - y^2
\]
La partie réelle est \( \dfrac{3}{4} - y^2 \)

Comme \( z^2 + z + 1 \) est réel, et \( Q = \dfrac{1}{z^2 + z + 1} \), \( Q \) est réel si \( z = -\dfrac{1}{2} + i y \).

Ainsi, les valeurs de \( z \) pour lesquelles les points \( z \), \( z^2 \) et \( z^4 \) sont alignés sont :

- Les nombres réels \( z \) (avec \( y = 0 \)).

- Les nombres complexes \( z = -\dfrac{1}{2} + i y \), avec \( y \in \mathbb{R} \).

\textbf{Conclusion :} Les nombres complexes \( z \) tels que \( z \in \mathbb{R} \) ou \( z = -\dfrac{1}{2} + i y \) avec \( y \in \mathbb{R} \) sont tels que les points \( z \), \( z^2 \) et \( z^4 \) sont alignés.

\end{solutionbox}

\end{document}

\end{document}